% =====================================================================
%  CSF-constrained single-image patch overfit — formal procedure
%
%  Corresponds exactly to:
%      python scripts/overfit.py --patch_mode csf --from_image \
%             --loss_fn cospgd --csf_threshold <tau> --steps <K>
%
%  Requires: amsmath, amssymb, and natbib-style \citep (swap for \cite
%  if using plain bibtex). Bibliography: csf_overfit_procedure.bib
%
%  PROVENANCE IS MARKED THROUGHOUT. Steps carrying a citation come from
%  the literature; steps marked [novel] have no external source and must
%  be defended in the text rather than cited. Sec.~11 tabulates this so
%  nothing is claimed by accident.
% =====================================================================

\subsection{Notation}

\begin{tabular}{ll}
$\mathbf{x}\in\mathbb{R}^{3\times H\times W}$ & input image, ImageNet-normalised \\
$\mathbf{x}^{01}=\mathrm{clip}(\mathbf{x}\odot\boldsymbol{\sigma}_n+\boldsymbol{\mu}_n,0,1)$ & the same image in $[0,1]$ RGB \\
$\mathbf{y}\in\{0,\dots,C-1\}\cup\{255\}^{H\times W}$ & ground-truth trainIds; $255$ is void \\
$f$ & frozen segmentation network, $f(\mathbf{x})\in\mathbb{R}^{C\times H'\times W'}$ \\
$\mathcal{U}$ & bilinear upsampling to $H\times W$ \\
$s,\; p=\lfloor Hs\rfloor,\; S$ & patch scale, rendered side, parameter resolution \\
$\mathbf{z}\in\mathbb{R}^{3\times S\times S}$ & the optimised parameter \\
$\mathcal{F},\mathcal{F}^{-1}$ & real 2-D DFT and inverse, forward-normalised \\
$\mathcal{K}$ & half-spectrum bin index set \\
$\tau$ & perceptual budget, in nominal JND (Sec.~\ref{sec:tau}) \\
\end{tabular}

% =====================================================================
\subsection{Step 1 — Observer model}

Contrast sensitivity is defined over cycles per degree of visual angle,
whereas an image spectrum is indexed in cycles per pixel. With physical
pixel size $\rho$ and viewing distance $d$, one pixel subtends
\begin{equation}
  \theta_{\mathrm{pix}} \;=\; 2\arctan\!\left(\frac{\rho}{2d}\right)\cdot\frac{180}{\pi}
  \qquad\text{degrees},
  \label{eq:cpd}
\end{equation}
so an image frequency $\nu$ in cycles per pixel corresponds to
$u=\nu/\theta_{\mathrm{pix}}$ cycles per degree
\citep[App.~D]{galinska2026csflow}. We adopt the Barten model
\citep{barten2003},
\begin{align}
  \mathrm{CSF}(u) &= \frac{M_{\mathrm{opt}}(u)}{k}
  \left[\frac{2}{T}\!\left(\frac{1}{X_0^{2}}+\frac{1}{X_{\max}^{2}}
  +\frac{u^{2}}{N_{\max}^{2}}\right)
  \!\left(\frac{1}{\eta\,\mathrm{p}\,E}
  +\frac{\Phi_0}{1-\exp\!\left(-(u/u_0)^{2}\right)}\right)\right]^{-1/2},
  \label{eq:barten}\\
  M_{\mathrm{opt}}(u) &= \exp\!\left(-2\pi^{2}\varsigma^{2}u^{2}\right),
\end{align}
with the parameterisation of \citet[App.~A]{galinska2026csflow}:
$\varsigma=0.5/60$, $k=3$, $T=0.1$, $X_0=60$, $X_{\max}=12$,
$N_{\max}=15$, $\eta=0.03$, $\mathrm{p}=1.2\times10^{6}$, $E=500$,
$\Phi_0=3\times10^{-8}$, $u_0=7$. Note $\mathrm{CSF}(0)=0$: a uniform
field carries no contrast. Where an orientation-dependent model is
required we substitute the Standard Spatial Observer with its oblique
effect \citep{watson2005sso}, under which diagonal carriers are less
visible than axis-aligned ones at equal radial frequency.

\paragraph{Scope of the observer assumption.}
Every visibility statement below is conditional on $(\rho,d)$ through
Eq.~\eqref{eq:cpd}. This is not an incidental detail:
\citet[Sec.~6]{galinska2026csflow} identify the assumption of a typical
screen and viewing distance as a \emph{conceptual limitation} of any
CSF-derived weighting, requiring calibration if optimal values are
wanted. We therefore treat $(\rho,d)$ as a reported experimental
parameter and sweep it, rather than fixing it once.

\paragraph{Why a pixel-space formulation is required.}
The correspondence between $\mathcal{F}$ bins and retinal spatial
frequency holds only when the perturbation lives in pixel space.
\citet[Sec.~5.1]{galinska2026csflow} make the same restriction for the
same reason: in latent-space models the coordinates of the perturbed
variable need not correspond to image spatial frequencies at all. The
patch here is optimised directly in $[0,1]$ RGB, so
Eq.~\eqref{eq:cpd} is well defined throughout.

% =====================================================================
\subsection{Step 2 — Per-frequency amplitude budget \hfill{\normalfont[novel]}}

A grating of Michelson contrast $c$ at frequency $u$ is at detection
threshold when $c\,\mathrm{CSF}(u)=1$ \citep{barten2003}. Inverting this
into an \emph{amplitude allowance} on the DFT is what turns a perception
model into an attack constraint. It does not appear in
\citet{galinska2026csflow}, who use $\mathrm{CSF}$ as a \emph{weight}
over generative timesteps; the function and its constants are taken from
that work, the inversion is not.

On the half-spectrum grid with radial frequency
$\nu(\mathbf{k})=\sqrt{\nu_y^2+\nu_x^2}$ in cycles per pixel,
\begin{equation}
  B(\mathbf{k}) \;=\;
  \min\!\left(\frac{\tau}{\kappa\,\mathrm{CSF}\!\big(\nu(\mathbf{k})/\theta_{\mathrm{pix}}\big)},\,A_{\max}\right)
  \cdot\mathbb{1}\!\left[\nu(\mathbf{k})\ \ge\ \nu_{\min}\right].
  \label{eq:budget}
\end{equation}

\paragraph{Contrast conversion $\kappa=4$.}
A cosine of peak amplitude $A$ on mean luminance $\bar{L}$ has Michelson
contrast $A/\bar{L}$ \citep{michelson1927}; under the forward-normalised
DFT that component has magnitude $A/2$. Taking $\bar{L}=\tfrac12$ for a
$[0,1]$ image gives $c=4\,|\mathcal{F}(\cdot)(\mathbf{k})|$. The mean is
assumed rather than measured locally, which \emph{under}-estimates
visibility on dark content and is absorbed into the calibration of
$\tau$.

\paragraph{Low-frequency cutoff $\nu_{\min}=m_c/S$.}
Since $\mathrm{CSF}\to0$ at DC, Eq.~\eqref{eq:budget} would otherwise
grant the near-DC bins the \emph{largest} allowance of any frequency.
The Barten rolloff is measured for full-field gratings, where the eye
adapts to a slow luminance ramp; inside a patch of side $S$ a component
below $1/S$ cycles per pixel completes no cycle and is a brightness
offset against the surround, which the patch border renders as a visible
edge. We require $m_c=2$ cycles across the patch. Empirically, omitting
this inverts the intended spectrum: measured radial power concentrates
at the lowest bin rather than the highest.

% =====================================================================
\subsection{Step 3 — Spectral reparameterisation \hfill{\normalfont[novel]}}

Rather than projecting onto the feasible set, we reparameterise so the
constraint holds by construction:
\begin{equation}
  Z=\mathcal{F}(\mathbf{z}),\qquad
  \hat{Z}(\mathbf{k})=B(\mathbf{k})\,
  \frac{Z(\mathbf{k})}{\sqrt{1+|Z(\mathbf{k})|^{2}}},\qquad
  \boldsymbol{\delta}_0=\mathcal{F}^{-1}(\hat{Z}),
  \label{eq:reparam}
\end{equation}
so $|\hat{Z}(\mathbf{k})|<B(\mathbf{k})$ for every $\mathbf{k}$ and any
$\mathbf{z}$. The scaling acts on the complex coefficient, preserving
phase exactly and avoiding differentiation of $\arg Z$ at the origin. A
hard projection would zero the gradient of every component resting on
its bound --- precisely where an adversarial component settles ---
whereas Eq.~\eqref{eq:reparam} is smooth everywhere and attains the
bound only in the limit.

% =====================================================================
\subsection{Step 4 — Visibility normalisation}

Detection of a compound stimulus is not decided by its most visible
component: the visual system pools evidence across spatial-frequency and
orientation channels, standardly modelled as a Minkowski norm with
exponent $\beta\approx3\text{--}4$
\citep{quick1974,robson1981,watson1979}. With $\mathcal{M}$ the set of
pixels actually composited,
\begin{equation}
  V_\beta(\boldsymbol{\delta})=
  \left[\sum_{\mathbf{k}\in\mathcal{K}}
  \Big(\kappa\,\big|\mathcal{F}(\boldsymbol{\delta}\odot\mathcal{M})(\mathbf{k})\big|\;
  \mathrm{CSF}\big(\nu(\mathbf{k})/\theta_{\mathrm{pix}}\big)\Big)^{\beta}
  \right]^{1/\beta},
  \label{eq:visibility}
\end{equation}
and the residual is rescaled to exactly the requested budget:
\begin{equation}
  \boldsymbol{\delta}_1=\frac{\tau}{V_\beta(\boldsymbol{\delta}_0)}\,\boldsymbol{\delta}_0 .
  \label{eq:scale}
\end{equation}
Shape and scale are deliberately separate: Eq.~\eqref{eq:reparam} bounds
each bin but says nothing about their sum, while Eq.~\eqref{eq:scale}
alone would place energy wherever the attack prefers. Equality rather
than inequality is used because the constraint is active at every step
and equality is differentiable where $\min(\cdot)$ is not.

% =====================================================================
\subsection{Step 5 — Dynamic-range fit \hfill{\normalfont[novel]}}

The composite must lie in $[0,1]$, but clipping is a hard nonlinearity
and generates broadband harmonics at exactly the frequencies where
$\mathrm{CSF}$ peaks, destroying the guarantee of
Eqs.~\eqref{eq:budget}--\eqref{eq:scale}. A \emph{uniform} rescale is
the only operation leaving the spectral shape invariant, since it scales
every Fourier coefficient equally and hence
$V_\beta(c\boldsymbol{\delta})=c\,V_\beta(\boldsymbol{\delta})$.
With per-pixel headroom
\begin{equation}
  a(u,v)=\begin{cases}
    1-\mathbf{r}(u,v), & \delta_1(u,v)>0\\
    \mathbf{r}(u,v),   & \text{otherwise,}
  \end{cases}
  \qquad
  c^{\star}=\mathrm{clip}\!\left(
  Q_q\!\left\{\tfrac{a(u,v)}{|\delta_1(u,v)|}\ :\ (u,v)\in\mathcal{M}\right\},0,1\right),
  \label{eq:fit}
\end{equation}
we set $\boldsymbol{\delta}=c^{\star}\boldsymbol{\delta}_1$, where $Q_q$
is the $q$-quantile ($q=10^{-3}$). The quantile rather than the minimum
admits a negligible fraction of clipped pixels, since a single saturated
pixel would otherwise force $c^{\star}=0$. Restricting the quantile to
$\mathcal{M}$ is essential: pixels outside the composited region cannot
affect the result yet would otherwise determine its scale.

% =====================================================================
\subsection{Step 6 — Reference and patch}

The base is the image content the patch replaces, at the resolved
placement $(t,\ell)$:
\begin{equation}
  \mathbf{r}=\mathrm{Resize}_{S\times S}
  \Big(\mathbf{x}^{01}\big[\,t:t{+}p,\ \ell:\ell{+}p\,\big]\Big),
  \qquad
  \mathbf{p}=\mathrm{clip}\big(\mathbf{r}+\boldsymbol{\delta},\,0,\,1\big).
  \label{eq:patch}
\end{equation}
$\mathbf{r}$ is treated as a constant ($\nabla_{\mathbf{r}}\equiv0$).
Because $\mathbf{r}$ equals the covered content, replacement by
$\mathbf{p}$ is identical to adding $\boldsymbol{\delta}$, so
$\tau\to0$ recovers the unperturbed image exactly and the
no-attack baseline is a no-op by construction.

% =====================================================================
\subsection{Step 7 — Application}

Following the localised-patch threat model of \citet{brown2017patch},
with optional silhouette mask $\mathcal{M}$ after \citet{tan2021lap},
\begin{align}
  \tilde{\mathbf{p}} &= \big(\mathrm{Resize}_{p\times p}(\mathbf{p})-\boldsymbol{\mu}_n\big)\oslash\boldsymbol{\sigma}_n, \\
  \mathbf{m} &= \mathrm{Pad}_{t,\ell}\big(\mathbb{1}_{p\times p}\odot\mathcal{M}\big), \\
  \mathbf{x}^{\mathrm{adv}} &= \mathbf{x}\odot(1-\mathbf{m})
  \;+\;\mathrm{Pad}_{t,\ell}(\tilde{\mathbf{p}})\odot\mathbf{m}.
  \label{eq:apply}
\end{align}

% =====================================================================
\subsection{Step 8 — Objective}

The scored support excludes both void labels and the patch footprint,
\begin{equation}
  \Omega=\{(u,v)\,:\,\mathbf{y}(u,v)\neq255\}\ \setminus\ \mathrm{supp}(\mathbf{m}),
\end{equation}
the latter to prevent the degenerate solution in which the patch itself
adopts the attacked appearance rather than influencing its surroundings
\citep{mirsky2023ipatch}. We use CosPGD \citep{agnihotri2024cospgd},
\begin{equation}
  \mathcal{L}=-\frac{1}{|\Omega|}\sum_{(u,v)\in\Omega}
  \underbrace{\Big[\cos\!\big(\mathrm{softmax}\,\mathcal{U}f(\mathbf{x}^{\mathrm{adv}})_{:,u,v},\,
  \mathbf{e}_{\mathbf{y}(u,v)}\big)\Big]_{\mathrm{sg}}}_{\text{alignment weight}}
  \cdot\ \mathrm{CE}\big(\mathcal{U}f(\mathbf{x}^{\mathrm{adv}})_{:,u,v},\,\mathbf{y}(u,v)\big).
  \label{eq:cospgd}
\end{equation}

\paragraph{Declared deviations.} \citet{agnihotri2024cospgd} do not
detach the cosine; $[\cdot]_{\mathrm{sg}}$ removes gradient flow through
the weight. They further use sign-SGD with $\varepsilon$-projection, as
a PGD variant \citep{madry2018pgd}; a patch is not an
$\varepsilon$-ball perturbation, so Eq.~\eqref{eq:cospgd} is optimised
with Adam \citep{kingma2015adam} on the raw gradient.

% =====================================================================
\subsection{Step 9 — Optimisation and evaluation}

\begin{equation}
  \mathbf{z}^{(0)}\sim\mathcal{N}(\mathbf{0},\mathbf{I}),\qquad
  \mathbf{z}^{(j+1)}=\mathrm{Adam}\big(\mathbf{z}^{(j)},\,
  \nabla_{\mathbf{z}}\mathcal{L}\big),\qquad j=0,\dots,K-1.
\end{equation}
The initialisation is Gaussian rather than zero because
Eq.~\eqref{eq:scale} is undefined at $\boldsymbol{\delta}_0=\mathbf{0}$;
Eqs.~\eqref{eq:reparam}--\eqref{eq:scale} are scale-invariant in
$\mathbf{z}$, so only the \emph{shape} of the initialisation matters.
$f$ is frozen throughout and $\nabla_{\phi}f$ is never formed. The
attack is reported as
\begin{equation}
  \Delta\mathrm{mIoU}_{\mathrm{remote}}
  =\mathrm{mIoU}\big(\mathcal{U}f(\mathbf{x}),\mathbf{y};\,\Omega\big)
  -\mathrm{mIoU}\big(\mathcal{U}f(\mathbf{x}^{\mathrm{adv}}),\mathbf{y};\,\Omega\big),
\end{equation}
the mean taken over classes present in $\mathbf{y}|_\Omega$ so both
terms share a denominator \citep{cordts2016cityscapes}. Spectral
placement of the learned residual is reported by radially averaged power
spectral density \citep{ruzanski2011rapsd}, and perceptual similarity to
$\mathbf{r}$ by LPIPS \citep{zhang2018lpips}, computed after training
and never entering Eq.~\eqref{eq:cospgd}.

% =====================================================================
\subsection{The budget parameter $\tau$}
\label{sec:tau}

\paragraph{Definition.} $\tau$ is the permitted Minkowski visibility of
the residual, in units of the Barten detection threshold under the
assumed geometry: by Eq.~\eqref{eq:scale} the composited perturbation
satisfies $V_\beta(\boldsymbol{\delta}_1)=\tau$ exactly, so $\tau=1$
nominally places the perturbation \emph{at} threshold and $\tau<1$
below it. It plays the role that $\alpha$ plays in
\citet{tan2021lap} --- the single knob trading attack strength against
perceptibility --- and the realism/strength curve over $\tau$ is the
deliverable of this family, not any single operating point.

\paragraph{$\tau$ acts once, not twice.} Although $\tau$ appears in both
Eq.~\eqref{eq:budget} and Eq.~\eqref{eq:scale}, its occurrence in the
budget is inert whenever the cap $A_{\max}$ is inactive: $B$ is then
proportional to $\tau$, hence so are $\boldsymbol{\delta}_0$ and
$V_\beta(\boldsymbol{\delta}_0)$, and the ratio in
Eq.~\eqref{eq:scale} cancels it. Measured at $S=128$ and the default
geometry, varying $\tau$ in Eq.~\eqref{eq:budget} alone over
$[0.05,1]$ leaves the residual bit-identical, and $A_{\max}=0.25$ binds
on no bin until $\tau\gtrsim4$. Eq.~\eqref{eq:budget} therefore sets the
\emph{shape} of the residual and Eq.~\eqref{eq:scale} its \emph{scale}.

\paragraph{Operating range.} Below the knee of
Eq.~\eqref{eq:fit}, realised visibility tracks $\tau$ exactly and the
residual scales linearly ($\tau=0.05,0.1,0.25,0.5$ give RMS
$0.0049,0.0098,0.0245,0.0490$). Above it the range fit binds, the
attained visibility saturates, and $\tau$ ceases to control anything ---
the knee is reference-dependent, since it is set by the headroom of
$\mathbf{r}$, and lay near $\tau\approx0.4$ for the references used
here. Values above the knee should not be reported as if they were
achieved.

\paragraph{Why $\tau$ is \emph{nominal}.} Three approximations stand
between $\tau$ and a psychophysical JND, and each biases it: the mean
luminance in $\kappa$ is assumed rather than measured; the Minkowski sum
in Eq.~\eqref{eq:visibility} runs over DFT bins rather than perceptual
channels; and Eq.~\eqref{eq:cpd} fixes a viewing geometry that
\citet[Sec.~6]{galinska2026csflow} identify as needing calibration.
$\tau$ should therefore be calibrated against rendered patches and
reported as a relative axis. A claim of imperceptibility requires a
study with human observers, which this procedure does not constitute.

\paragraph{Not to be confused with the interpolation parameter of
\citet{galinska2026csflow}.} Their $\alpha$ in
$w_{\mathrm{final}}=\alpha w_{\mathrm{CSFlow}}+(1-\alpha)w_{\mathrm{base}}$
blends a perception-driven weighting with a baseline one over generative
timesteps. It is not a perceptual budget and does not bound any
quantity; the two parameters share no role beyond both being scalars
derived from the same CSF.

\paragraph{A residual bounded by $\mathrm{CSF}$ need not be
statistically inconspicuous.} Natural image spectra follow an
approximate power law, so energy placed near Nyquist is atypical of the
data even where the eye is least sensitive
\citep{falck2025fourier,dieleman2024spectral}. Invisibility to a human
observer and indistinguishability from natural image statistics are
distinct properties, and only the former is constrained here; the RAPSD
of Sec.~9 is what makes the latter measurable.

% =====================================================================
\subsection{Provenance}

\begin{tabular}{lll}
\textbf{Step} & \textbf{Element} & \textbf{Source} \\
\hline
1 & cycles/pixel $\to$ cycles/degree & \citep[App.~D]{galinska2026csflow} \\
1 & Barten CSF; parameter values & \citep{barten2003}; params \citep[App.~A]{galinska2026csflow} \\
1 & Standard Spatial Observer, oblique effect & \citep{watson2005sso} \\
1 & pixel-space requirement; geometry caveat & \citep[Sec.~5.1, 6]{galinska2026csflow} \\
2 & \emph{inversion of CSF into a budget} & \textbf{novel} \\
2 & Michelson contrast, $\kappa=4$ & \citep{michelson1927} \\
2 & \emph{low-frequency cutoff for a patch} & \textbf{novel} \\
3 & \emph{smooth spectral reparameterisation} & \textbf{novel} \\
4 & Minkowski / probability summation & \citep{quick1974,robson1981,watson1979} \\
5 & \emph{uniform-rescale range fit} & \textbf{novel} \\
6 & \emph{covered content as residual base} & \textbf{novel} (cf.\ \citep{tan2021lap}) \\
7 & localised patch application & \citep{brown2017patch} \\
7 & silhouette masking & \citep{tan2021lap} \\
8 & footprint exclusion & \citep{mirsky2023ipatch} \\
8 & CosPGD objective & \citep{agnihotri2024cospgd} \\
8 & PGD lineage; declared deviation & \citep{madry2018pgd} \\
9 & Adam & \citep{kingma2015adam} \\
9 & mIoU protocol, Cityscapes & \citep{cordts2016cityscapes} \\
9 & RAPSD & \citep{ruzanski2011rapsd} \\
9 & LPIPS, evaluation only & \citep{zhang2018lpips} \\
10 & natural-image spectral statistics & \citep{falck2025fourier,dieleman2024spectral} \\
\end{tabular}

\medskip
\noindent
Five elements carry no external source and are contributions of this
work: the budget inversion (Eq.~\ref{eq:budget}), the patch-specific
low-frequency cutoff, the smooth spectral reparameterisation
(Eq.~\ref{eq:reparam}), the uniform-rescale range fit
(Eq.~\ref{eq:fit}), and the use of the covered content as the residual
base (Eq.~\ref{eq:patch}). From \citet{galinska2026csflow} we take a
function, its constants, a unit conversion and two scoping caveats ---
not a method: their contribution is a weighting over flow-matching
timesteps, which has no analogue here.
